\begin{proof}[Proof of \cref{obj:lem:positive-endpoint}]
Fix \(0<\beta\le1\), and put
\[
c_{\mathrm{env}}:=\frac{\pi^6}{4}.
\]
We first establish positivity, unit mass, and the square-integral estimate with constant \(c_{\mathrm{env}}\), and then prove the moment estimate with a positive constant \(C_{\mathrm{bias}}\). The common constant will be
\[
C:=\max\{C_{\mathrm{bias}},c_{\mathrm{env}}\}.
\]
Since \(L^{-2\beta}\) and \(L^2\) are nonnegative, enlarging either constant to \(C\) preserves its respective bound.
% lean: CausalSmith.Stat.RdTruesideNoiseFrontier.positive_endpoint

\begin{enumerate}
\item Fix an integer \(L\ge1\). Define the polynomial quotient and its normalizer by
\[
F_L(x):=
\begin{cases}
\dfrac{1-T_L(1-2x)}{2x},&x\ne0,\\[6pt]
L^2,&x=0,
\end{cases}
\qquad x\in\mathbb R,
\qquad
c_{\mathrm{norm},L}:=\int_0^1 F_L(t)^3\,dt.
\]
Indeed, \(T_L(1)=1\), so the numerator is divisible by \(x\). The standard Chebyshev derivative identity gives
\[
T_L'(1)=L^2,
\qquad
2xF_L(x)=1-T_L(1-2x),
\qquad
F_L(0)=T_L'(1)=L^2,
\]
where the last equality also follows by differentiating the preceding polynomial identity at zero.

The normalization in \cref{obj:synth_4} implies, for \(0<x\le1\),
\[
T_L(1-2x)
=
\cos\!\bigl(L\arccos(1-2x)\bigr)
\le1.
\]
Thus \(F_L(x)\ge0\) on \((0,1]\), and continuity extends this inequality to zero. Consequently,
\[
F_L(x)\ge0\quad(0\le x\le1),
\qquad
c_{\mathrm{norm},L}>0.
\]
Strict positivity of the integral follows from continuity, nonnegativity, and \(F_L(0)^3=L^6>0\). By \cref{obj:def:endpoint-kernel},
\[
K_L(x)=c_{\mathrm{norm},L}^{-1}F_L(x)^3,
\]
and hence
\[
K_L(x)\ge0\quad(0\le x\le1),
\qquad
\int_0^1K_L(x)\,dx
=
c_{\mathrm{norm},L}^{-1}c_{\mathrm{norm},L}
=1.
\]
% lean: CausalSmith.Stat.RdTruesideNoiseFrontier.endpointKernel_integral_one

\item We next obtain a quantitative lower bound for the normalizer. For \(0<x\le1\), define
\[
\varphi_x:=\arcsin\sqrt{x}.
\]
Then
\[
0\le\varphi_x\le\frac{\pi}{2},
\qquad
\sin\varphi_x=\sqrt{x}.
\]
The polynomial identity above, the Chebyshev identity in \cref{obj:synth_4}, and the double-angle formula give
\[
2\sin^2\varphi_x\,F_L(\sin^2\varphi_x)
=
1-\cos(2L\varphi_x)
=
2\sin^2(L\varphi_x).
\]
In particular,
\[
xF_L(x)=\sin^2(L\varphi_x).
\]

Define the endpoint-window width by
\[
\delta_{\mathrm{win},L}:=\frac{1}{16L^2}.
\]
For \(0<x\le\delta_{\mathrm{win},L}\), we have
\[
L\sqrt{x}\le\frac14.
\]
Using \(\sin v\ge(2/\pi)v\) for \(0\le v\le\pi/2\) at \(v=\varphi_x\), we obtain
\[
\frac{2}{\pi}L\varphi_x
\le L\sin\varphi_x
=L\sqrt{x}
\le\frac14.
\]
Thus \(0\le L\varphi_x\le\pi/2\), so the same sine inequality applies at \(L\varphi_x\). Together with \(\sin\varphi_x\le\varphi_x\), it yields
\[
\frac{2}{\pi}L\sin\varphi_x
\le\frac{2}{\pi}L\varphi_x
\le\sin(L\varphi_x).
\]
Squaring and using \(xF_L(x)=\sin^2(L\varphi_x)\), then dividing by \(x>0\), gives
\[
F_L(x)\ge\left(\frac{2}{\pi}\right)^2L^2.
\]
At \(x=0\), this inequality follows from \(F_L(0)=L^2\) and \(2/\pi\le1\). Therefore it holds throughout \([0,\delta_{\mathrm{win},L}]\). Since \(F_L^3\) is nonnegative on \([0,1]\),
\[
\begin{aligned}
c_{\mathrm{norm},L}
&\ge
\int_0^{\delta_{\mathrm{win},L}}F_L(t)^3\,dt\\
&\ge
\int_0^{\delta_{\mathrm{win},L}}
\left(\left(\frac{2}{\pi}\right)^2L^2\right)^3\,dt\\
&=
\frac{1}{16L^2}\frac{64L^6}{\pi^6}
=
\frac{4}{\pi^6}L^4.
\end{aligned}
\]
% lean: CausalSmith.Stat.RdTruesideNoiseFrontier.endpointNormalizer_lower_bound

\item To obtain the uniform kernel bound, we use
\[
|\sin(L\varphi_x)|\le L|\sin\varphi_x|.
\]
This follows by induction on the nonnegative integer \(L\). The zero case is immediate, and the sine addition formula gives the induction step:
\[
\begin{aligned}
|\sin((L+1)\varphi_x)|
&\le
|\sin(L\varphi_x)|\,|\cos\varphi_x|
+
|\cos(L\varphi_x)|\,|\sin\varphi_x|\\
&\le
|\sin(L\varphi_x)|+|\sin\varphi_x|\\
&\le
(L+1)|\sin\varphi_x|.
\end{aligned}
\]
For \(0<x\le1\), squaring this bound gives
\[
xF_L(x)=\sin^2(L\varphi_x)
\le L^2\sin^2\varphi_x
=L^2x,
\]
so \(F_L(x)\le L^2\). At zero, equality holds. Hence
\[
F_L(x)^3\le L^6
\le c_{\mathrm{norm},L}\,c_{\mathrm{env}}L^2
\qquad(0\le x\le1),
\]
where the second inequality is the normalizer lower bound multiplied by \(c_{\mathrm{env}}L^2\). Dividing by the positive normalizer proves
\[
K_L(x)\le c_{\mathrm{env}}L^2
\qquad(0\le x\le1).
\]
% lean: CausalSmith.Stat.RdTruesideNoiseFrontier.endpointKernel_le_index_sq

\item Multiplying the uniform bound by the nonnegative kernel and integrating gives
\[
\begin{aligned}
\int_0^1K_L(x)^2\,dx
&\le
\int_0^1 c_{\mathrm{env}}L^2K_L(x)\,dx\\
&=
c_{\mathrm{env}}L^2\int_0^1K_L(x)\,dx
=
c_{\mathrm{env}}L^2.
\end{aligned}
\]
% lean: CausalSmith.Stat.RdTruesideNoiseFrontier.endpointKernel_square_integral_bound

\item For the moment estimate, we also need a tail bound. For \(0<x\le1\), the Chebyshev identity gives
\[
T_L(1-2x)
=
\cos\!\bigl(L\arccos(1-2x)\bigr)
\ge-1.
\]
Consequently,
\[
2xF_L(x)=1-T_L(1-2x)\le2,
\qquad
F_L(x)\le x^{-1}.
\]
Cubing this inequality and using the normalizer lower bound yields
\[
F_L(x)^3
\le x^{-3}
\le c_{\mathrm{norm},L}\,c_{\mathrm{env}}L^{-4}x^{-3}.
\]
Division by \(c_{\mathrm{norm},L}>0\) proves
\[
K_L(x)\le c_{\mathrm{env}}L^{-4}x^{-3}
\qquad(0<x\le1).
\]
% lean: CausalSmith.Stat.RdTruesideNoiseFrontier.endpointKernel_tail_bound

\item Define the splitting point by
\[
\delta_L:=L^{-2},
\qquad
0<\delta_L\le1.
\]
On \([0,\delta_L]\), monotonicity of \(x^\beta\), nonnegativity of \(K_L\), and the uniform kernel bound give
\[
x^\beta K_L(x)
\le
\delta_L^\beta c_{\mathrm{env}}L^2.
\]
Therefore,
\[
\begin{aligned}
\int_0^{\delta_L}x^\beta K_L(x)\,dx
&\le
\int_0^{\delta_L}
\delta_L^\beta c_{\mathrm{env}}L^2\,dx\\
&=
c_{\mathrm{env}}\delta_L^\beta,
\end{aligned}
\]
because \(\delta_LL^2=1\).
% lean: CausalSmith.Stat.RdTruesideNoiseFrontier.endpointKernel_moment_bound

\item On \([\delta_L,1]\), the tail bound gives
\[
x^\beta K_L(x)
\le
c_{\mathrm{env}}\delta_L^2x^{\beta-3}.
\]
Since \(\delta_L>0\) and \(0<\beta\le1\), integration of this power yields
\[
\begin{aligned}
\int_{\delta_L}^1x^\beta K_L(x)\,dx
&\le
c_{\mathrm{env}}\delta_L^2
\int_{\delta_L}^1x^{\beta-3}\,dx\\
&=
c_{\mathrm{env}}\delta_L^2
\frac{\delta_L^{\beta-2}-1}{2-\beta}.
\end{aligned}
\]
Here \(2-\beta\ge1\) and \(\delta_L^{\beta-2}\ge0\), so
\[
\delta_L^{\beta-2}-1
\le
(2-\beta)\delta_L^{\beta-2}.
\]
It follows that
\[
\int_{\delta_L}^1x^\beta K_L(x)\,dx
\le
c_{\mathrm{env}}\delta_L^2\delta_L^{\beta-2}
=
c_{\mathrm{env}}\delta_L^\beta.
\]
These identities also cover \(L=1\): then \(\delta_L=1\), and the integral over \([\delta_L,1]\) is zero.
% lean: CausalSmith.Stat.RdTruesideNoiseFrontier.endpointKernel_moment_bound

\item The function \(x^\beta K_L(x)\) is continuous on \([0,1]\), so its integral splits at \(\delta_L\). Combining the two bounds gives
\[
\begin{aligned}
M_{L,\beta}
&=
\int_0^{\delta_L}x^\beta K_L(x)\,dx
+
\int_{\delta_L}^1x^\beta K_L(x)\,dx\\
&\le
2c_{\mathrm{env}}\delta_L^\beta
=
2c_{\mathrm{env}}L^{-2\beta}.
\end{aligned}
\]
Thus we may take
\[
C_{\mathrm{bias}}:=2c_{\mathrm{env}}>0,
\qquad
C:=\max\{C_{\mathrm{bias}},c_{\mathrm{env}}\}>0.
\]
The moment and square-integral estimates then become
\[
M_{L,\beta}\le CL^{-2\beta},
\qquad
\int_0^1K_L(x)^2\,dx\le CL^2.
\]
Together with positivity and unit mass, these establish all the assertions for every integer \(L\ge1\).
% lean: CausalSmith.Stat.RdTruesideNoiseFrontier.positive_endpoint
\end{enumerate}
\end{proof}
